\subsection{Exact nearest-neighbor crossing sums}
\label{sec:area_law_edge_crossing_budget}

\begin{definition}[On-site and nearest-neighbor supports]
    \label{def:area_law_nearest_neighbor_support}
    \lean{TNLean.PEPS.AreaLaw.nearestNeighborSupport}
    \leanok
    \uses{def:area_law_domain}
    Let $\Lambda\subset\mathbb Z^2$ be finite and let $E_\Lambda$ be its
    unordered nearest-neighbor edges. Index supports by the disjoint union
    $I=\Lambda\sqcup E_\Lambda$, assigning $S_v=\{v\}$ to a site and
    $S_{\{x,y\}}=\{x,y\}$ to an edge. For $A\subseteq\Lambda$, put
    $C(A)=\{i\in I:S_i\cap A\ne\varnothing,\ S_i\setminus A\ne\varnothing\}$.
\end{definition}

\begin{theorem}[One inside endpoint]
    \label{thm:area_law_crossing_inside_endpoint}
    \lean{
        TNLean.PEPS.AreaLaw.not_mem_crossingTerms_singleton,
        TNLean.PEPS.AreaLaw.pair_inter_eq_singleton_of_crossing,
        TNLean.PEPS.AreaLaw.card_pair_insideConfig_of_crossing
    }
    \leanok
    A singleton cannot meet both sides of a cut in any finite set.
    If $\{x,y\}$ meets both sides of $A$, then
    $\{x,y\}\cap A=\{v\}$ for exactly one endpoint $v$.
    For an alphabet of cardinality $q\ge0$, the set of configurations on
    this intersection has cardinality $q$.
\end{theorem}
\begin{proof}
    \leanok
    If $x\in A$, the crossing condition forces $y\notin A$;
    otherwise it forces $y\in A$. A function on a singleton is determined
    by its value at the unique point, including when the alphabet is empty.
\end{proof}

\begin{theorem}[Crossing supports and unordered edges]
    \label{thm:area_law_crossing_support_count}
    \lean{
        TNLean.PEPS.AreaLaw.inl_not_mem_crossingTerms,
        TNLean.PEPS.AreaLaw.inr_mem_crossingTerms_iff,
        TNLean.PEPS.AreaLaw.card_nearestNeighbor_insideConfig,
        TNLean.PEPS.AreaLaw.card_crossingTerms_nearestNeighbor
    }
    \leanok
    \uses{def:area_law_nearest_neighbor_support,def:area_law_hamiltonian}
    For every $A\subseteq\Lambda$, no on-site index belongs to $C(A)$,
    and an edge index belongs to $C(A)$ exactly when the edge lies in
    $\partial_\Lambda A$. Consequently $|C(A)|=|\partial_\Lambda A|$.
    For every $i\in C(A)$ and every integer $q\ge0$, the configuration
    set on $S_i\cap A$ has cardinality $q$.
\end{theorem}
\begin{proof}
    \leanok
    \uses{thm:area_law_crossing_inside_endpoint}
    Choose any linear order on $\Lambda$ and write each edge in increasing
    endpoint order. Forgetting this order maps crossing edge indices to
    unordered boundary edges. Equality of unordered pairs either identifies
    both endpoints or swaps them; the swap contradicts their strict order.
    Every boundary edge has an increasing ordering, so this map is bijective.
    The configuration count follows from the unique inside endpoint.
\end{proof}

\begin{theorem}[Exact weighted crossing sum]
    \label{thm:area_law_nearest_neighbor_budget}
    \lean{TNLean.PEPS.AreaLaw.nearestNeighbor_crossing_budget}
    \leanok
    \uses{def:area_law_nearest_neighbor_support,def:area_law_hamiltonian}
    Let $m,q\ge0$ be integers, $A_j\subseteq\Lambda$ for $0\le j<m$,
    and let $a_j,J\in\mathbb R$. If $d_{ij}$ denotes the cardinality of
    the configuration set on $S_i\cap A_j$ for an alphabet of size $q$, then
    \begin{align}
        \sum_{j=0}^{m-1}a_j^2\sum_{i\in C(A_j)}d_{ij}^2J
        &=q^2J\sum_{j=0}^{m-1}a_j^2|\partial_\Lambda A_j|.
        \label{eq:area_law_exact_crossing_sum}
    \end{align}
    This is the combinatorial identity used in the nearest-neighbor
    matrix-unit estimate of \cite[Section~4]{OpenAI2026PolynomialPEPS}.
\end{theorem}
\begin{proof}
    \leanok
    \uses{thm:area_law_crossing_support_count}
    Every summand in the inner sum has $d_{ij}=q$, and there are exactly
    $|\partial_\Lambda A_j|$ such summands. Substitute these identities
    and factor out $q^2J$.
\end{proof}
